\chapter{Background Knowledge}
\label{chap:background}

% PLAN: written alongside the clusters (W2--W6), ~15--18 pages.
% Intro paragraph: this chapter introduces the concepts needed to follow
% Chapter~\ref{chap:literature-review}. Notation is fixed here and reused
% everywhere else.

\section{Supervised Learning Setting and Notation}
\label{sec:bg-notation}

% (W2) Input space, label, classifier, sensitive attribute, cost function
% between an instance and its counterfactual, set of feasible actions.
% Keep it short; this is the reference for all later formulas.

\section{Kolmogorov-Arnold Networks}
\label{sec:bg-kan}

\subsection{The Kolmogorov-Arnold Representation Theorem}
\label{subsec:bg-ka-theorem}

% (W3) Statement of the theorem \cite{kolmogorov1957representation, arnold1957functions};
% why it was long considered irrelevant for learning \cite{girosi1989representation}.

\subsection{Architecture and Training}
\label{subsec:bg-kan-architecture}

% (W3) Learnable univariate edge functions (B-splines + residual basis),
% grid extension, KAN layers vs. MLP layers \cite{liu2025kan};
% caveats from fair comparisons \cite{yu2024kan}. Figure: KAN vs. MLP schematic.

\subsection{Interpretability: Pruning and Symbolic Regression}
\label{subsec:bg-kan-interpretability}

% (W3) Sparsification, pruning, symbolic formula extraction
% \cite{liu2025kan, liu2024kan}; relation to additive models
% \cite{agarwal2021neural}.

\section{Algorithmic Fairness}
\label{sec:bg-fairness}

\subsection{Sources of Bias}
\label{subsec:bg-bias-sources}

% (W5) Taxonomy of bias sources \cite{mehrabi2021survey, barocas2023fairness}.

\subsection{Group Fairness Criteria}
\label{subsec:bg-group-fairness}

% (W5) Demographic parity, equal opportunity, equalized odds, calibration
% \cite{hardt2016equality}; impossibility results
% \cite{kleinberg2017inherent, chouldechova2017fair}.

\subsection{Individual and Causal Fairness}
\label{subsec:bg-individual-causal-fairness}

% (W5) Fairness through awareness \cite{dwork2012fairness}; counterfactual
% fairness \cite{kusner2017counterfactual}; causal reasoning about
% discrimination \cite{kilbertus2017avoiding}.

\subsection{Bias Mitigation Strategies}
\label{subsec:bg-mitigation}

% (W5) Pre-, in- and post-processing \cite{caton2024fairness};
% fair representations \cite{zemel2013learning}.

\section{Causal Inference Foundations}
\label{sec:bg-causality}

\subsection{Structural Causal Models}
\label{subsec:bg-scm}

% (W6) SCMs, causal graphs, structural equations, exogenous noise
% \cite{pearl2009causality, peters2017elements}.

\subsection{Interventions and Counterfactuals}
\label{subsec:bg-counterfactuals}

% (W6) Ladder of causation; do-operator; abduction--action--prediction.
% Figure: a small SCM with an intervention (reuse it in Section~\ref{sec:bg-recourse}).

\section{Counterfactual Explanations and Algorithmic Recourse}
\label{sec:bg-recourse}

\subsection{Counterfactual Explanations}
\label{subsec:bg-ce}

% (W4) Optimisation formulation \cite{wachter2018counterfactual}; desiderata.

\subsection{From Counterfactual Explanations to Recourse}
\label{subsec:bg-ce-to-recourse}

% (W4) Actionability and cost \cite{ustun2019actionable}; recourse as minimal
% interventions instead of nearest counterfactuals \cite{karimi2021algorithmic}.

\subsection{Evaluation Criteria}
\label{subsec:bg-recourse-metrics}

% (W4) Validity, proximity, sparsity, plausibility, diversity, cost, runtime
% \cite{verma2024counterfactual, guidotti2024counterfactual}.

\section{Summary}
\label{sec:bg-summary}

% Half a page: how the four building blocks connect to the thesis.
